\uuid{KxoX}
\titre{Limites de suites}
\auteur{bodin}
\organisation{exo7}
\licence{CC BY-NC-SA 4.0}
\datecreate{2020-01-21}
\niveau{L1}
\module{Analyse}
\chapitre{Suite}
\sousChapitre{Convergence}
\difficulte{}
\theme{Réels}
\exo7qcmid{arnaud.arnaud.023}

\contenu{
\begin{question}
\qclassification{121.01}
Quelles sont les assertions vraies ?
\begin{answers}
    \good{Si $u_n= n\sin(\frac 1 n)$ alors $(u_n)$ tend vers $1$.}
    \bad{Si $u_n=\ln (\ln(n))$ alors $(u_n)$ a une limite finie.}
    \bad{$u_n = \frac{(\ln n)^2}{\sqrt n}$ alors $(u_n)$ tend vers $+\infty$.}
    \bad{$u_n = 1+\frac12 +\frac14 +\frac18 + \cdots +\frac{1}{2^n}$ alors $(u_n)$ diverge.}
    \good{$u_n = \sin(n)$, il existe une sous-suite de $(u_n)$ convergente.}
\end{answers}
\end{question}
}
